
\begin{tcolorbox}[
    enhanced,
    colback=white,
    colframe=black,
    arc=5mm,
    boxrule=1pt,
    left=2mm,
    right=2mm,
    top=1mm,
    bottom=1mm
]

\textbf{Task.} Assess how well the candidate response answers the question from a user's perspective. This rubric is \emph{orthogonal} to source faithfulness; here you evaluate relevance, completeness, correctness, and clarity.

\textbf{Rubric (integer 1--5).}\\
\textbf{5 Excellent.} Fully addresses every sub-question, factually correct, clearly written, appropriately detailed.\\
\textbf{4 Good.} Addresses the main point accurately; may miss a minor nuance or be marginally less clear than ideal.\\
\textbf{3 Partial.} Handles some sub-questions but omits or mishandles others, or answers a slightly different question than asked.\\
\textbf{2 Poor.} Contains substantial factual errors or largely fails to address the question.\\
\textbf{1 Inadequate.} Irrelevant, off-topic, or uninformative.\\

\textbf{Procedure.}\\
1. Decompose the question into its explicit and implicit sub-questions.\\
2. Check whether each sub-question is answered, not answered, or wrongly answered.\\
3. Consider style and clarity only as a tie-breaker between adjacent rubric levels.\\
4. Output the final line as \texttt{Rating: <integer 1-5>}.

\textbf{Demonstration.}\\
\emph{Passage (context only):} Modern Singapore was founded in 1819 by Raffles as a British trading post.\\
\emph{Question:} Tell me about Singapore's founding and its current major industries.\\
\emph{Response:} Singapore was founded by Raffles in 1819. Its primary industry is currently agriculture and fishing.\\
\emph{Sub-question coverage:} founding — correct; current major industries — incorrect (modern Singapore is known for finance, biomedical, electronics, not agriculture).\\
\emph{Rating: 3}

\textbf{Now evaluate the real instance.}\\

\textbf{Passage (context):}\\
\{DOCUMENT\}\\

\textbf{Question:}\\
\{QUESTION\}\\

\textbf{Response:}\\
\{RESPONSE\}\\

\end{tcolorbox}
